\documentclass[10pt,a4paper]{article}
\usepackage[margin=2cm]{geometry}
\usepackage{needspace}
\usepackage{fontspec,booktabs,longtable,array,xcolor,hyperref}
\setmainfont{TeX Gyre Pagella}
\setsansfont{TeX Gyre Heros}
\definecolor{reportblue}{HTML}{244D65}
\hypersetup{colorlinks=true,urlcolor=reportblue,linkcolor=reportblue}
\setlength{\parindent}{0pt}
\setlength{\parskip}{5pt}
\renewcommand{\arraystretch}{1.18}
\setlength{\tabcolsep}{5pt}
\emergencystretch=2em
\urlstyle{same}
\pagestyle{plain}
\title{\sffamily Jev on MoCa: Agreement with Human Causal and Moral Judgments}
\author{}
\date{September 30, 2026}
\begin{document}
\begin{center}{\sffamily\LARGE Jev on MoCa: Agreement with Human Causal and Moral Judgments\par}\vspace{8pt}{\small September 30, 2026}\end{center}\vspace{6pt}
\section*{Abstract}
Across 206 English MoCa stories, Jev’s three-class agreement with the aggregated judgments of 25 human respondents per item is 43.69\%. Agreement is 39.58\% for causal items and 53.23\% for moral items. Jev distinguishes clear human yes/no judgments but rarely places human-ambiguous causal stories in the ambiguous interval.
\section{Dataset}
The dataset includes 144 causal stories and 62 moral-permissibility stories, each with 25 binary human judgments. The classes are Yes, No and Ambiguous, using the 0.4/0.6 boundaries; Jev’s third class is derived from its binary-choice probability. Model snapshot: typesafe/jev-1.13-20260917.

Factor names follow the source’s expert annotations, and factors may overlap within a story. Human response proportions and Jev’s structured-choice probabilities have different meanings; this report describes their agreement and numerical differences on these stories.

\section{Results}
\Needspace{8\baselineskip}
\subsection{Overall results}
Agreement is judged against the dataset’s aggregated human labels, without treating majority judgments as moral truth. Uniform three-class chance is 33.33\%; the overall majority-class baseline is 36.41\%. AUROC uses only human-unambiguous items.

{\small
\begin{longtable}{@{}>{\raggedright\arraybackslash}p{\dimexpr 0.2400\textwidth-2\tabcolsep\relax}>{\raggedright\arraybackslash}p{\dimexpr 0.1520\textwidth-2\tabcolsep\relax}>{\raggedright\arraybackslash}p{\dimexpr 0.1520\textwidth-2\tabcolsep\relax}>{\raggedright\arraybackslash}p{\dimexpr 0.1520\textwidth-2\tabcolsep\relax}>{\raggedright\arraybackslash}p{\dimexpr 0.1520\textwidth-2\tabcolsep\relax}>{\raggedright\arraybackslash}p{\dimexpr 0.1520\textwidth-2\tabcolsep\relax}@{}}
\toprule
\textbf{Subset} & \textbf{Items} & \textbf{Agreement} & \textbf{95\% interval} & \textbf{AUROC} & \textbf{Probability MAE}\\
\midrule
\endfirsthead
\toprule
\textbf{Subset} & \textbf{Items} & \textbf{Agreement} & \textbf{95\% interval} & \textbf{AUROC} & \textbf{Probability MAE}\\
\midrule
\endhead
All & 206 & 43.69\% & 36.89\%–50.49\% & 0.704 & 0.308\\
Causal & 144 & 39.58\% & 31.94\%–47.92\% & 0.748 & 0.333\\
Moral & 62 & 53.23\% & 40.32\%–66.13\% & 0.852 & 0.249\\
\bottomrule
\end{longtable}
}
\Needspace{8\baselineskip}
\subsection{Judgment distributions}
{\small
\begin{longtable}{@{}>{\raggedright\arraybackslash}p{\dimexpr 0.2400\textwidth-2\tabcolsep\relax}>{\raggedright\arraybackslash}p{\dimexpr 0.1900\textwidth-2\tabcolsep\relax}>{\raggedright\arraybackslash}p{\dimexpr 0.1900\textwidth-2\tabcolsep\relax}>{\raggedright\arraybackslash}p{\dimexpr 0.1900\textwidth-2\tabcolsep\relax}>{\raggedright\arraybackslash}p{\dimexpr 0.1900\textwidth-2\tabcolsep\relax}@{}}
\toprule
\textbf{Subset} & \textbf{Source} & \textbf{Yes} & \textbf{No} & \textbf{Ambiguous}\\
\midrule
\endfirsthead
\toprule
\textbf{Subset} & \textbf{Source} & \textbf{Yes} & \textbf{No} & \textbf{Ambiguous}\\
\midrule
\endhead
Causal & Human & 48 & 50 & 46\\
Causal & Jev & 98 & 35 & 11\\
Moral & Human & 23 & 10 & 29\\
Moral & Jev & 22 & 21 & 19\\
\bottomrule
\end{longtable}
}
Only 1 of 46 human-ambiguous causal items is also classified as ambiguous by Jev; the corresponding count is 12 of 29 moral items. Jev gives substantially more affirmative causal judgments, showing that lower three-class agreement involves distribution and boundary differences.

\Needspace{8\baselineskip}
\subsection{Mean differences across factor groups}
{\small
\begin{longtable}{@{}>{\raggedright\arraybackslash}p{\dimexpr 0.3200\textwidth-2\tabcolsep\relax}>{\raggedright\arraybackslash}p{\dimexpr 0.2267\textwidth-2\tabcolsep\relax}>{\raggedright\arraybackslash}p{\dimexpr 0.2267\textwidth-2\tabcolsep\relax}>{\raggedright\arraybackslash}p{\dimexpr 0.2267\textwidth-2\tabcolsep\relax}@{}}
\toprule
\textbf{Group transition} & \textbf{Group sizes} & \textbf{Human difference} & \textbf{Jev difference}\\
\midrule
\endfirsthead
\toprule
\textbf{Group transition} & \textbf{Group sizes} & \textbf{Human difference} & \textbf{Jev difference}\\
\midrule
\endhead
Action to omission & 80 / 32 & +0.028 & -0.179\\
Aware to unaware & 23 / 20 & -0.009 & +0.195\\
Conjunctive to disjunctive & 66 / 32 & -0.030 & -0.250\\
Abnormal to normal & 44 / 45 & -0.228 & -0.294\\
Other to self beneficiary & 26 / 22 & +0.101 & -0.015\\
Accidental to instrumental & 18 / 30 & -0.118 & -0.067\\
Avoidable to inevitable & 22 / 26 & +0.131 & +0.136\\
Impersonal to personal force & 24 / 24 & -0.058 & -0.067\\
\bottomrule
\end{longtable}
}
Differences describe changes in mean affirmative share/probability from the first group to the second. Factors overlap and story content changes across groups; these are descriptive associations, not isolated causal effects.

\Needspace{8\baselineskip}
\subsection{Call costs}
{\small
\begin{longtable}{@{}>{\raggedright\arraybackslash}p{\dimexpr 0.2400\textwidth-2\tabcolsep\relax}>{\raggedright\arraybackslash}p{\dimexpr 0.3800\textwidth-2\tabcolsep\relax}>{\raggedright\arraybackslash}p{\dimexpr 0.3800\textwidth-2\tabcolsep\relax}@{}}
\toprule
\textbf{Input tokens} & \textbf{Output tokens} & \textbf{Reported cost (USD)}\\
\midrule
\endfirsthead
\toprule
\textbf{Input tokens} & \textbf{Output tokens} & \textbf{Reported cost (USD)}\\
\midrule
\endhead
109,940 & 6,798 & 0.00461748\\
\bottomrule
\end{longtable}
}
All 206 requests succeeded on the first attempt, with no unknown charges or retries.

\Needspace{8\baselineskip}
\section{Conclusion}
Jev shows partial agreement with aggregated human judgments, with higher agreement on moral than causal items. Its ability to distinguish clear human yes/no judgments does not translate into comparable identification of human ambiguity. Some factor-group mean directions match and others differ; the findings describe judgments on the same stories.
\Needspace{5\baselineskip}
\section*{References}
\begin{enumerate}
\item Nie et al. (2023). MoCa: Measuring Human-Language Model Alignment on Causal and Moral Judgment Tasks. NeurIPS. \url{https://arxiv.org/abs/2310.19677}
\item MoCa official repository, commit 1b61a20294247480d64675ceb19751ef4e1e878f. \url{https://github.com/cicl-stanford/moca}
\end{enumerate}
\end{document}
